\Subhead{3.7.1}{Тепловой шум. Классический и квантовый варианты формулы Найквиста}

Резистор в равновесии описывают последовательной случайной ЭДС. Для сопротивления $R$ в классической области односторонняя плотность в герцах
\begin{equation}
 G_U(f)=4k_BTR,\qquad G_I(f)=4k_BT/R.
\end{equation}
Пример проверки: RC-цепь с шумом резистора имеет $\langle U_C^2\rangle=\int_0^\infty4k_BTR/[1+(2\pi fRC)^2]\dd f=k_BT/C$. Энергия $C\langle U_C^2\rangle/2=k_BT/2$ соответствует равнораспределению. Для пассивного линейного двухполюсника в тепловом равновесии обобщения различаются:
\begin{equation}\label{eq:nyquistZY}
 G_U=4k_BT\operatorname{Re}Z(f),\qquad
 G_I=4k_BT\operatorname{Re}Y(f),\qquad Y=1/Z.
\end{equation}
Действительно, эквивалентные источники Тевенина и Нортона связаны $I=U/Z$, поэтому $G_I=G_U/|Z|^2$ и $\operatorname{Re}Y=\operatorname{Re}Z/|Z|^2$. Это не замена $R$ на $\operatorname{Re}Z$ в знаменателе токовой формулы.
% Самостоятельный вывод из эквивалентности источников и преобразования спектра.

Квантовая \emph{симметризованная} односторонняя формула:
\begin{equation}
 G_U^{\rm sym}(f)=4\operatorname{Re}Z(f)\,
 \frac{hf}{2}\coth\frac{hf}{2k_BT}.
\end{equation}
Здесь содержится $hf/2$ нулевых флуктуаций. Доступная согласованной нагрузке тепловая мощность в узкой полосе $\Delta f$ определяется заселённостью $\bar n=[\exp(hf/k_BT)-1]^{-1}$:
\begin{equation}
 P_{\rm av}=hf\bar n\,\Delta f\longrightarrow k_BT\Delta f.
\end{equation}
Нельзя переносить нулевой член симметризованного спектра в формулу извлекаемой мощности. Выбор спектра зависит от способа измерения. При согласовании классическая мощность не зависит от $R$: рост ЭДС компенсируется делением напряжения и величиной нагрузки.
